\section{Background}
\label{sec:background}

\subsection{Machine-consumed instruction files}
\todo{Define AGENTS.md, skills, IDE rules; contrast with README/CONTRIBUTING (P5 descriptive baseline). Cite practitioner docs.}

\subsection{Reference integrity and truth decay}
Longitudinal panels label verifiable references as \textsc{Verified} or \textsc{Missing} at each commit.
Prior work in this laboratory distinguishes:
\begin{itemize}
  \item \textbf{Born-stale} references: \textsc{Missing} without a prior \textsc{Verified} observation.
  \item \textbf{Post-verification decay}: transition from \textsc{Verified} to \textsc{Missing}.
\end{itemize}
RQ3 explicitly warns against merging these prevalence rates.

\subsection{Agent maintenance attribution}
Gate P4 validates whether commits reflect genuine agent-assisted instruction maintenance.
Human gold validation on $N=200$ commits yields precision 0.958 and F1 0.9617 for the positive class \texttt{true\_agent\_maintenance} (\path{exports/truth_pilot/p4_validation.md}).

\subsection{Controlled agent impact experiments}
RQ5 manipulates instruction-file content at a pinned commit while holding repository state, task prompt, and test command fixed.
The redesigned protocol adds condition~C (no instruction file) to separate instruction \emph{presence} from referential \emph{truth}.
